\section{Introduction}
Recent advances in large language models and planning architectures have enabled agentic systems that reason over extended horizons, invoke external tools, modify internal state, and coordinate with other agents. Benchmarks such as ALFWorld, WebArena, and SWE-bench reflect this shift, evaluating agents on increasingly complex, multi-step tasks. As a result, contemporary agent evaluation is dominated by task-level success metrics---accuracy, completion rate, reward, or test pass rate---which serve as effective indicators of capability.

However, as agent autonomy and complexity increase, performance metrics alone are insufficient to answer a growing class of questions: \emph{Why did an agent take a particular action?} \emph{Can its decision process be reconstructed after execution?} \emph{Did the agent’s behavior change over time, and if so, how explicitly?} \emph{In multi-agent settings, which agent was responsible for which decisions?} These questions are not about whether an agent succeeded, but whether its reasoning process is \emph{inspectable, attributable, and stable under analysis}.

Existing work on interpretability, logging, and agent frameworks provides partial tools for examining internal behavior, but these efforts are typically ad hoc, architecture-specific, or tightly coupled to particular training methods. Moreover, they rarely offer a standardized way to compare reasoning structure across agents, tasks, or benchmarks. As a result, two agents with similar performance may differ substantially in how transparent or reconstructable their reasoning processes are, yet this distinction is not captured by prevailing evaluation practices.

In this paper, we introduce \emph{reasoning security} as an evaluation dimension orthogonal to performance. Reasoning security concerns the extent to which an agent’s decisions can be causally reconstructed and analyzed post hoc, using available traces, logs, or artifacts, without imposing constraints on how the agent is built or trained. Importantly, reasoning security does \emph{not} assess behavioral desirability, safety, or alignment, nor does it prescribe how agents should act. Instead, it focuses on structural properties of decision processes that determine whether reasoning can be audited, compared, and understood after execution.

To operationalize this concept, we propose the \emph{Reasoning Security Score (RSS)}, a composite diagnostic measure derived from five decomposable dimensions: traceability, lineage integrity, adaptation visibility, attribution fidelity, and counterfactual stability.\footnote{RSS is computed post hoc from logged decision traces and does not affect model training, inference, or benchmark execution; it is optional and intended solely to contextualize agent behavior alongside performance metrics.} RSS is computed post hoc from decision records and does not alter training objectives, inference procedures, or benchmark execution. It is designed to be architecture-agnostic, minimally invasive, and applicable to both single-agent and multi-agent systems.

The contributions of this work are threefold. First, we formalize reasoning security as a distinct evaluation dimension for agentic systems and clarify how it differs from performance, interpretability, and control-oriented frameworks. Second, we introduce a concrete instrumentation model and metric suite that enables consistent measurement of reasoning security across benchmarks. Third, we demonstrate how RSS can be applied to common agent evaluation settings, showing that decomposed reasoning security metrics reveal meaningful structural differences between agents that are not apparent from performance alone.

By positioning RSS as complementary evaluation infrastructure rather than a replacement for existing metrics, this work aims to support deeper analysis of agent behavior as autonomy and complexity continue to scale.
